\usepackage[utf8]{inputenc} %кодировка
\usepackage{cmap} % поиск в PDF - перемещен перед fontenc
\usepackage[T1,T2A]{fontenc} % тоже кодировка
\usepackage[english,russian]{babel} %языки

\usepackage{multirow}
\usepackage{tabularx,booktabs}
\usepackage{longtable}

% Используем ulem только один раз с нужными опциями
\usepackage[normalem]{ulem}
\usepackage{graphicx, float} %пакет для единого оформления всех плавающих объектов (избегаем повторяющихся команд в документе)

% Добавляем пакет textcomp для символов \textperthousand и \textmu
\usepackage{textcomp}

\DeclareGraphicsExtensions{.pdf,.png,.jpg,.eps}%форматы 
\usepackage{titlesec}
\setcounter{secnumdepth}{4}
\titleformat{\paragraph}
{\normalfont\normalsize\bfseries}{\theparagraph}{1em}{}
\titlespacing*{\paragraph}
{0pt}{3.25ex plus 1ex minus .2ex}{1.5ex plus .2ex}

\graphicspath{{pictures/}} % выбираем папку, в которую сохраняем все рисунки, чтобы не было хаоса файлов

\usepackage{amsmath,amssymb,amsfonts,amsthm,mathtools} %математические формулы и символы

\usepackage[a4paper,left=30mm,right=15mm,top=20mm,bottom=20mm]{geometry} % устанавливает поля документа

\parindent=4ex %красная строка
\parskip=3mm %расстояние между параграфами
\usepackage{indentfirst} %делать отступ в начале параграфа

% Убедимся, что hyperref загружается последним из пакетов
% для предотвращения конфликтов
\usepackage{hyperref} % добавление ссылок

\def\hmath$#1${\texorpdfstring{{\rmfamily\textit{#1}}}{#1}}
%настройка подписей плавающих объектов

\makeindex %нумерация 

\usepackage{array,caption} %картинки, подписи, таблицы
%\usepackage{endfloat} - для вывода картинок со списком в конце файла

\captionsetup[table]{skip=4pt,singlelinecheck=off}

% Используем только subcaption, убираем конфликтующий subfigure
\usepackage[labelformat=simple]{subcaption} %для subfigure
\renewcommand\thesubfigure{(\alph{subfigure})}

\usepackage[export]{adjustbox} %чтобы влево-вправо картинки ставить

\usepackage{afterpage,placeins} % для барьеров 

\usepackage{wrapfig} %добавление wrapfig

\usepackage[nottoc]{tocbibind} %подключает в содержание список лит-ры

\usepackage{multicol}

\usepackage{setspace}

\usepackage{gensymb}

\usepackage{listings}

\usepackage{mathtext} % русские буквы в формулах
\usepackage{icomma} % "Умная" запятая

%% Шрифты
\usepackage{euscript}	 % Шрифт Евклид
\usepackage{mathrsfs} % Красивый матшрифт

%% Свои команды
\DeclareMathOperator{\sgn}{\mathop{sgn}}

%% Перенос знаков в формулах (по Львовскому)
\newcommand*{\hm}[1]{#1\nobreak\discretionary{}
	{\hbox{$\mathsurround=0pt #1$}}{}}

\usepackage{verbatim}

\usepackage{xcolor}
\usepackage{colortbl}

\usepackage{siunitx} % Required for alignment
% Удален конфликтующий пакет subfigure
\usepackage{rotating}

%New colors defined below
\definecolor{codegreen}{rgb}{0,0.6,0}
\definecolor{codegray}{rgb}{0.5,0.5,0.5}
\definecolor{codepurple}{rgb}{0.58,0,0.82}
\definecolor{backcolour}{rgb}{0.95,0.95,0.92}

%Code listing style named "mystyle"
\lstdefinestyle{mystyle}{
  backgroundcolor=\color{backcolour}, commentstyle=\color{codegreen},
  keywordstyle=\color{magenta},
  numberstyle=\tiny\color{codegray},
  stringstyle=\color{codepurple},
  basicstyle=\ttfamily\footnotesize,
  breakatwhitespace=false,         
  breaklines=true,                 
  captionpos=b,                    
  keepspaces=true,                 
  numbers=left,                    
  numbersep=5pt,                  
  showspaces=false,                
  showstringspaces=false,
  showtabs=false,                  
  tabsize=2
}

%"mystyle" code listing set
\lstset{style=mystyle}